\begin{tutorial}{Молекулярная биология}

\subsection*{1 подгруппа}

Заметим, что если мы начинаем набирать слово с позиции $s$, то не существует двух различных слов из словаря
$\overline{a_s a_{s+1} \ldots a_{s+i}}$
и
$\overline{a_s a_{s+1} \ldots a_{s+j}}$,
так как одно из них было бы префиксом другого, что запрещено по условию.

Следовательно, из каждой позиции строки однозначно определяется слово словаря, которое может начинаться в этой позиции (если оно вообще существует).

Предподсчитаем для каждой позиции $i$:
\begin{itemize}
    \item существует ли слово словаря, начинающееся в $i$,
    \item если существует --- позицию его конца.
\end{itemize}

Это можно сделать прямой проверкой за $O(SL)$, где $S$ --- суммарная длина слов словаря, $L$ --- длина белка.

Теперь обработка одного запроса $[l, r]$ происходит так:
\begin{itemize}
    \item начинаем с позиции $l$,
    \item переходим к концу найденного слова,
    \item продолжаем, пока текущая позиция не станет равной $r+1$.
\end{itemize}

Если мы перепрыгнули за $r+1$, разбиение невозможно.

Таким образом, сложность: $O(SL + LM)$

\subsection*{2 подгруппа}

Предподсчёт из предыдущего решения слишком дорогой.

Построим бор по всем словам словаря. Тогда мы можем проходить по белку посимвольно.

Алгоритм:
\begin{itemize}
    \item идём по буквам строки,
    \item двигаемся по бору,
    \item если пришли в терминальную вершину, значит найдено слово,
    \item переносим указатель в корень бора и увеличиваем ответ.
\end{itemize}

Так как ни одно слово не является префиксом другого, терминальность однозначно определяет завершение слова.

Предобработка занимает $O(S)$, а каждый проход по фрагменту --- $O(L)$.

Итоговая сложность: $O(S + LM)$.

\subsection*{3 подгруппа}

Используем хеширование.

Сохраним в \texttt{set} или \texttt{unordered\_set} хеши всех префиксов слов словаря. Также посчитаем префиксные хеши строки белка.

Для позиции $s$ можно бинарным поиском найти длину слова, начинающегося в $s$, сравнивая хеши за $O(\log S)$ и проверять, есть ли он в \texttt{set}.

Заметим, что из каждой позиции может начинаться не более одного слова.

Построим ориентированное ребро $i \to j+1$, если из позиции $i$ начинается слово, заканчивающееся в позиции $j$.

Так как из каждой вершины выходит не более одного ребра, получаем лес.

Тогда:
\begin{itemize}
    \item разбиение существует тогда и только тогда, когда вершина $l$ лежит в поддереве вершины $r+1$,
    \item это можно проверить через времена входа и выхода ($tin$, $tout$),
    \item ответ равен разности глубин.
\end{itemize}

Сложность:
\[
O(S \log S + L \log^2 S + Q)
\]
или
\[
O(S + L \log S + Q),
\]
в зависимости от структуры хранения хешей.


\subsection*{Полное решение}

Заметим, что структура переходов полностью совпадает с автоматом Ахо--Корасик.

Построим автомат по словарю за
\[
O(SA),
\]
где $A$ --- размер алфавита.

Суперсуффиксные ссылки (output links) соответствуют тем самым переходам из решения 3 подгруппы.
После каждого перехода по суперсуффиксной ссылке мы добавляем одно ребро в лес.

Тогда:
\begin{itemize}
    \item при проходе по белку мы получаем все окончания слов,
    \item строим граф переходов,
    \item далее отвечаем на запросы так же, как в 3 подгруппе.
\end{itemize}
Тогда суперсуффиксных переходов $O(L)$, так как из каждого индекса $i$ выходит не более одного ребра.

Итоговая сложность:
\[
O(SA + (L + S) + Q).
\]



\end{tutorial}
